\chapter{序論}
\label{chap:intro}

本章は，卒業論文・修士論文における「序論（はじめに）」の書き方を実例として示すものである．
序論は読者が最初に目にする章であり，論文全体の印象を大きく左右する．
一般に，\textbf{研究背景}・\textbf{問題提起}・\textbf{研究目的}・\textbf{論文の構成}の順に記述するとよい\cite{kisaragi}．
本章自体がその構成例になっている．

%% ===========================================================
\section{研究の背景}
\label{sec:background}
まず，研究対象とする分野の現状や社会的な背景を述べ，読者を研究テーマへ導く．
本テンプレートは，\upLaTeX\footnote{\LaTeX を日本語（特に縦書きや多書体）に対応させた処理系である．}
を用いて論文を執筆するための環境を提供する．
組版システムとしての\TeX は文献\cite{lamport}や文献\cite{tex}をはじめ多くの文献で解説されており，
学術論文の作成において広く利用されている．

背景の記述では，専門外の読者にも伝わるように，一般的な事柄から具体的な対象へと
順に話を進めることが望ましい．略語を用いる場合は初出時に正式名称を併記する
（例：自然言語処理 (Natural Language Processing; NLP)）．

%% ===========================================================
\section{既存研究の課題}
\label{sec:problem}
次に，既存の手法や状況にどのような課題があるのかを示す．
ここでの問題提起が，後の研究目的と対応している必要がある．
課題を述べる際は，単なる感想ではなく，先行研究の引用\cite{tex,kisaragi}に基づいて
客観的に論じることが重要である．

%% ===========================================================
\section{本研究の目的}
\label{sec:purpose}
問題提起を受けて，本研究が何を目指すのかを明確に述べる．
目的が曖昧であったり，問題提起と対応していなかったりする論文は再考を要する．
本テンプレートの目的は，論文の体裁を整える手間を最小化し，
執筆者が内容の検討に集中できる環境を提供することである．

%% ===========================================================
\section{本論文の構成}
\label{sec:structure}
最後に，以降の章で何を述べるのかを示す．
読者はこの記述を地図として論文を読み進める．

\begin{description}
    \item[第\ref{chap:intro}章\ 序論] 研究の背景・課題・目的と，本論文の構成を述べる（本章）．
    \item[第\ref{chap:writing}章\ 論文執筆の作法] 論文全体の構成や，わかりやすい文章を書くための作法をまとめる．
    \item[第\ref{chap:latex}章\ 文章作成のための\LaTeX 記法] 見出し・参照・箇条書き・引用など，文章を組み立てる基本的な記法を示す．
    \item[第\ref{chap:float}章\ 図・表・数式・コードの記法] 図の載せ方や表・数式・ソースコードなど，論文で多用する記法をパターン豊富に示す．
    \item[第\ref{chap:conclusion}章\ 結論] 本論文のまとめと今後の課題を述べる．
\end{description}

\begin{itembox}[l]{序論に含めるべき要素}
    \begin{enumerate}
        \item 研究背景：分野の現状や社会的背景
        \item 問題提起：既存手法・状況の課題
        \item 研究目的：課題に対して本研究が目指すこと
        \item 論文の構成：以降の章の概要
    \end{enumerate}
\end{itembox}
